%%%%%%%%%%%%%%%%%%%%%%%%%%%%%%%%%%%%%%%%%%%%%%%%%%%%%%%%%%%%%%%%%%%%%%
% Lớp 10 | Phương pháp tọa độ trong mặt phẳng (𝑂𝑥𝑦) | Tọa độ của véctơ
% Dạng 0H9?1-5: Phân tích một véctơ theo 2 véctơ không cùng phương
% Số câu: 5  (sinh tự động bởi tach_cau.py; từ nay thầy sửa trực tiếp file này)
%%%%%%%%%%%%%%%%%%%%%%%%%%%%%%%%%%%%%%%%%%%%%%%%%%%%%%%%%%%%%%%%%%%%%%

%%%%%%%%%%%%%%%% ĐÚNG/SAI %%%%%%%%%%%%%%%%

% nguồn: MaHoaNH_DungSai_Lop10
\begin{ex}%[0H9N1-5]%[ma:2d76e291]
	Cho hai véc-tơ $ \overrightarrow{u}=(2;-3) $ và $ \overrightarrow{v}=(1;4) $. Biết $\overrightarrow{OM}= \overrightarrow{u}-2\overrightarrow{v}$. Xác định tính đúng sai của  các mệnh đề sau?
	 \choiceTF
		 {$\overrightarrow{OM}=(0;11)$}
		{\True  $M(0;-11)$}
		{\True Điểm $M$ tồn tại duy  nhất}
		{Điểm $M$ không tồn tại} 
	 \loigiai{
		Ta có $ \overrightarrow{u}-2\overrightarrow{v} =(0;-11)$. 
		\begin{itemchoice}
			\itemch $\overrightarrow{OM}=(0,-11)$. 
			 \itemch $M(0,-11)$.
			\itemch $M(0,-11)$  tồn tại duy nhất.
			\itemch $M(0,-11)$:  tồn tại.
		\end{itemchoice}
	 }
\end{ex}

% nguồn: MaHoaNH_DungSai_Lop10
\begin{ex}%[0H9N1-5]%[ma:2b7f6885]
	Trong mặt phẳng tọa độ $Oxy$, cho điểm $A(4;0)$, $B(4;2)$ và $C(0;2)$. Xét tính đúng sai của các mệnh đề sau 
	 \choiceTF
		{\True  $\overrightarrow{OA}=(4;0)$, $\overrightarrow{OB}=(4;2)$, $\overrightarrow{OC}=(0;2)$}
		{\True $O, A, C$  không thẳng hàng}
		{$\overrightarrow{OA}+3\overrightarrow{OC}= \overrightarrow{OB}$} 
		{\True  Tồn tại số thực $k$ sao cho $\overrightarrow{OA}+k\overrightarrow{OC}= 2\overrightarrow{OB}$}
	\loigiai{ 
		\begin{itemchoice}
			\itemch  $\overrightarrow{OA}=(4;0)$, $\overrightarrow{OB}=(4;2)$, $\overrightarrow{OC}=(0;2)$.
			\itemch  Có
			$\overrightarrow{OA}=(4;0)$  và $\overrightarrow{OC}=(0;2)$ nên hai véc-tơ không cùng phương.\\
			Do đó, $O, A, C$ không thẳng  hàng.
			 \itemch
			 $\overrightarrow{OA}+3\overrightarrow{OC}=(4;6)$.\\
			$\overrightarrow{OB}=(4;2)$ 
			\itemch $\overrightarrow{OA}=(4;0)$,  $k\overrightarrow{OC}=(0;2k)$, $2\overrightarrow{OB}=(4;2)$.\\
			Với $k=1$ thì $\overrightarrow{OA}+k\overrightarrow{OC}= 2\overrightarrow{OB}$. 
		\end{itemchoice}
	} 
\end{ex}

% nguồn: MaHoaNH_DungSai_Lop10
\begin{ex}%[0H9H1-5]%[ma:4f223e0b]
	Trong  mặt phẳng tọa độ $Oxy$, cho các vectơ $\overrightarrow{a}=(2;-2)$, $\overrightarrow{b}=(4; 1)$ và $\overrightarrow{c}=(0;-1)$. Khi đó
	 \choiceTF
		{\True $2\overrightarrow{a}-\overrightarrow{b}-3\overrightarrow{c}=(0;-2)$} 
		{\True  $\overrightarrow{e}=(1;-1)$ cùng phương, cùng hướng với $\overrightarrow{a}$}
		{$\overrightarrow{f}=\left(-1;-\dfrac{1}{4}\right)$  cùng phương, cùng hướng với $\overrightarrow{b}$}
		{\True $\overrightarrow{a}=\dfrac{1}{2}  \overrightarrow{b}+\dfrac{5}{2} \overrightarrow{c}$}
	\loigiai{ 
		\begin{itemchoice}
			\itemch Đúng. Ta có $\heva{&2 \overrightarrow{a}=(4;-4) \\&-\overrightarrow{b}=(-4;-1) \\&-3 \overrightarrow{c}=(0; 3)}  \Rightarrow \overrightarrow{d}=2 \overrightarrow{a}-\overrightarrow{b}-3 \overrightarrow{c}=(0;-2)$.
			\itemch Đúng. Ta có $\overrightarrow{a}=(2;-2)=2 \overrightarrow{e}$ nên $\overrightarrow{a}$, $\overrightarrow{e}$ là hai vectơ cùng phương với nhau, hơn nữa  chúng cùng hướng với nhau vì $\overrightarrow{a}=k \overrightarrow{e}$, $k=2>0$.
			\itemch Sai. Ta có $\overrightarrow{b}=(4; 1)=-4 \overrightarrow{f}$, tức  là $\overrightarrow{b}=k \overrightarrow{f}$, $k=-4<0$ nên $\overrightarrow{b}$ và $\overrightarrow{f}$ là hai vectơ cùng phương, ngược hướng với nhau.
			\itemch Đúng. Gọi $m$, $n$ là các số thỏa mãn $\overrightarrow{a}=m \overrightarrow{b}+n \overrightarrow{c}$  ($\overrightarrow{b}$, $\overrightarrow{c}$ không cùng phương).\\
			Khi đó $\heva{&2=m \cdot  4+n \cdot 0 \\&-2=m \cdot 1+n \cdot(-1)} \Leftrightarrow\heva{&m=\dfrac{1}{2} \\&n=\dfrac{5}{2}}$. Vậy $\overrightarrow{a}=\dfrac{1}{2} \overrightarrow{b}+\dfrac{5}{2} \overrightarrow{c}$.
		\end{itemchoice}
	} 
\end{ex}

% nguồn: MaHoaNH_DungSai_Lop10
\begin{ex}%[0H9H1-5]%[ma:a730b85c]
	Trong mặt phẳng tọa độ $Oxy$, cho  các điểm $A(-4; 1)$, $B(2; 4)$, $C(2;-2)$. Khi đó
	\choiceTF 
		{Tọa độ điểm $D$ sao cho $C$ là trọng tâm tam giác $ABD$  là $D(8;11)$}
		{\True Tọa độ điểm $E$ thuộc trục hoành  sao cho $A$, $B$, $E$ thẳng hàng là $E(-6; 0)$}
		{\True $\overrightarrow{BC}=(0;-6)$, $\overrightarrow{AC}=(6;-3)$} 
		{Tọa độ $F$ thỏa mãn $\overrightarrow{A F}=\overrightarrow{B C}-2 \overrightarrow{A C}+2 \overrightarrow{C F}$ là  $F(20;5)$}
	 \loigiai{
		\begin{itemchoice}
			\itemch Sai. $C$ là trọng  tâm tam giác $ABD$\\
			$\Leftrightarrow \heva{&x_C=\dfrac{x_A+x_B+x_D}{3}\\&y_C=\dfrac{y_A+y_B+y_D}{3}}\Leftrightarrow  \heva{&2=\dfrac{-4+2+x_D}{3}\\&-2=\dfrac{1+4+y_D}{3}}\Leftrightarrow \heva{&x_D=8\\&y_D=-11}$.\\
			Vậy  $D(8;-11)$.
			\itemch Đúng. Gọi $E(x;  0) \in Ox\Rightarrow \overrightarrow{A E}=(x+4;-1)$, $\overrightarrow{A B}=(6; 3)$.\\
			Ba điểm $A$, $B$, $E$ thẳng hàng $\Leftrightarrow \overrightarrow{A E}$ cùng phương $\overrightarrow{A B}\\ \Leftrightarrow \dfrac{x+4}{6}=\dfrac{-1}{3}$ $\Leftrightarrow  x+4=-2 \Leftrightarrow x=-6$.\\
			Vậy $E(-6;  0)$.
			\itemch Đúng. Ta có $\overrightarrow{B C}=(0;-6)$, $\overrightarrow{A C}=(6;-3)$. 
			\itemch Sai. Gọi $F(x; y)$.  Ta có $\overrightarrow{A F}=(x+4; y-1)$, $2\overrightarrow{CF}=(2x-4;2y+4)$.\\
			Suy ra $\overrightarrow{B C}-2 \overrightarrow{A  C}+2 \overrightarrow{C F}=(2 x-16; 2 y+4)$.\\
			Ta có $\overrightarrow{A F}=\overrightarrow{B C}-2 \overrightarrow{A C}+2 \overrightarrow{C F} \Leftrightarrow\heva{&x+4=2 x-16  \\&y-1=2 y+4} \Leftrightarrow\heva{&x=20 \\&y=-5.}$\\
			Vậy  $F(20;-5)$.
		\end{itemchoice}
	} 
\end{ex}


%%%%%%%%%%%%%%%% TRẢ LỜI NGẮN %%%%%%%%%%%%%%%%

% nguồn: MaHoaNH_TraLoiNgan_Lop10
\begin{bt}%[0H9H1-5]%[ma:a7b889e6]
	Trong mặt phẳng toạ độ $Oxy$,  cho $\overrightarrow{a}=(m; 2)$, $\overrightarrow{b}=(-3; n)$ và $\overrightarrow{c}=(-2 m; 7)$. Tìm $m+n$ biết $\overrightarrow{c}=\overrightarrow{a}+\overrightarrow{b}$.
	 \shortans{6}
	\loigiai{ 
		Ta  có
		$$ 
		 \overrightarrow{c}=\overrightarrow{a}+\overrightarrow{b}
		 \Leftrightarrow
		\heva{&-2 m=m-3 \\ &7=2+n} 
		 \Leftrightarrow
		\heva{&m=1 \\  &n=5.}
		$$ 
		Vậy $m+n=6$. 
	} 
\end{bt}
